% Real points of E: y^2 = 4x^3 - 4x.  Two components -- a compact oval over [e1,e2] and an
% unbounded branch over [e3, infty) meeting at O at infinity.  The three roots together
% with O form the 2-torsion E[2] = {O, e1, e2, e3} = (Z/2)^2.
\begin{tikzpicture}[>=Latex, xscale=1.28, yscale=0.6]
  \draw[->, black!50, line width=0.5pt] (-1.6,0)--(2.5,0) node[right, font=\small] {$x$};
  \draw[->, black!50, line width=0.5pt] (0,-5.7)--(0,6.0) node[above, font=\small] {$y$};
  \draw[formteal!62!black, line width=1.3pt, smooth, samples=60, domain=-1:0]
    plot (\x,{ 2*sqrt(max(\x*\x*\x-\x,0))});
  \draw[formteal!62!black, line width=1.3pt, smooth, samples=60, domain=-1:0]
    plot (\x,{-2*sqrt(max(\x*\x*\x-\x,0))});
  \draw[formteal!62!black, line width=1.3pt, smooth, samples=80, domain=1:2.12]
    plot (\x,{ 2*sqrt(max(\x*\x*\x-\x,0))});
  \draw[formteal!62!black, line width=1.3pt, smooth, samples=80, domain=1:2.12]
    plot (\x,{-2*sqrt(max(\x*\x*\x-\x,0))});
  \foreach \r in {-1,0,1}{\fill[formrust!85!black] (\r,0) circle (2.5pt);}
  \node[below left=-1pt, font=\small, formrust!70!black] at (-1,0) {$e_1$};
  \node[below right=-1pt, font=\small, formrust!70!black] at (0,0) {$e_2$};
  \node[below right=-1pt, font=\small, formrust!70!black] at (1,0) {$e_3$};
  \node[formteal!55!black, font=\small] at (-0.5,1.65) {oval};
  \node[formteal!55!black, font=\small] at (1.95,2.0) {branch};
  \node[formslate, font=\small, anchor=south] at (2.05,5.2) {$\to O$};
  \node[formslate, font=\small, anchor=north] at (2.05,-5.0) {$\to O$};
\end{tikzpicture}
